\section{Introduction}\label{sec:intro}

Mixed models with subject-level random effects are routinely fitted by
Hamiltonian Monte Carlo (HMC) and its adaptive variant, the No-U-Turn
sampler \citep[NUTS;][]{hoffman2014}, in Stan \citep{carpenter2017}, brms
\citep{burkner2017}, rstanarm and NumPyro \citep{phan2019}. The standard
choice, and the default in brms and rstanarm, is the non-centered
parameterization $b = \Bstd L^{\top}$, with $\Bstd$ standard normal and $L$
a Cholesky factor of the random-effect covariance. It removes the funnel
between the random effects and their scale \citep{neal2003,betancourt2015}.

It leaves a second, simpler problem in place. Suppose the fixed-effects
design contains a column that is constant within subject: the intercept,
a treatment arm, sex. Adding a constant $c$ to every subject's random
intercept and subtracting $c$ from the fixed intercept leaves every fitted
value unchanged, so only the priors pin the split, at scale
$\sigma_0/\sqrt{N}$ for $N$ subjects when the fixed intercept's prior is
diffuse. The same holds for every subject-constant covariate, whose
coefficient trades off against a covariate-weighted combination of the
random intercepts, and, with a random slope on time, for the fixed time
slope and the mean random slope. In HMC's sampling coordinates these directions are spread across
all $N$ subjects, the posterior has a ridge along them, and the sampler
can mix slowly on the intercept, the subject-constant coefficients and the
time slope: the coefficients an analyst reports. In a joint model of the
\texttt{pbc2} data, for instance, the posterior correlation between the
fixed intercept and the mean random intercept was $-0.97$
(Section~\ref{sec:degeneracy}).

The problem is well understood for Gibbs samplers. Hierarchical centering
\citep{gelfand1995} removes it at the price of the funnel, and the choice
between centered and non-centered coordinates has a large literature
\citep{papaspiliopoulos2003,papaspiliopoulos2007,yu2011}. Other approaches
remove the directions: sum-to-zero constraints, whose effect on Gibbs and
NUTS is quantified by \citet{zanella2021}; restricted spatial regression
\citep{reich2006,hodges2010}; and, exactly, the transformation of
\citet{vines1996}, which integrates the grand-mean directions out.
\citet{yang2021} use the covariate-weighted direction to analyze Gibbs
samplers, and \citet{browne2009} name the problem, suggesting that
predictors be made orthogonal to the cluster indicators.

This paper keeps the absorbable directions and makes them sampling
coordinates. We construct an orthonormal basis for every direction the
fixed effects can absorb, in every random-effect column, and apply one
fixed rotation to the standardized random effects so that these
directions become a small block $U$ of $k \times q$ coordinates, with $k$
directions and $q$ random-effect columns. A single dense block of the HMC
mass matrix over $(\beta, U)$ then largely whitens the ridge. The rotation
acts on the subject index and the covariance factor on the column index,
so they commute: the construction is exact for every value of the
random-effect covariance, with no constraint, no correction and no change
to any reported quantity.

The motivating application is the Bayesian joint model of a longitudinal
marker and a time-to-event outcome
\citep{wulfsohn1997,henderson2000,tsiatis2004,rizopoulos2012}, as fitted
with NUTS by the R package \texttt{jmjax}. There the survival submodel adds
a condition on which directions remain absorbable, and a second question
arises. On the association parameter $\alpha$, the primary estimand, NUTS
gave 3 to 8 times the effective sample size (ESS) per second of the
Metropolis-within-Gibbs sampler of JMbayes2 \citep{jmbayes2} without any
reparameterization. We explain most of the per-draw part of that gap, on
two datasets where we measure the relevant quantity, with an ideal
two-block sampler, whose lag-1 autocorrelation for $\alpha$ equals
the share of $\alpha$'s posterior variance that the other parameters
explain \citep{liu1994}.

Section~\ref{sec:model} sets out the model and the degeneracy.
Section~\ref{sec:absorbable} characterizes the absorbable directions for
any design, including one where a natural column-by-column check misses
the slope direction, and the condition a survival submodel imposes.
Section~\ref{sec:rotation} gives the rotation, which costs $O(Nkq)$ per
gradient, and an analysis of the ridge that explains why a diagonal metric
cannot handle it, why one dense block can, and where a fixed metric is
weakest. Section~\ref{sec:blocked} treats $\alpha$.
Sections~\ref{sec:simulation}--\ref{sec:realdata} give evidence from
simulation, two mixed models and two joint-model datasets, including a
comparison with JMbayes2, and Section~\ref{sec:discussion} discusses
related work and limits. \suppref{sec:repro} lists the script that
reproduces each result.
